
\subsection*{投稿時第7節}
\Comment{R1.31}
「これは前の節の議論の特別な場合であり、ずっと短くできる。……説明が長すぎると、何を理解すべきか、前の節と何が違うのかという焦点がぼやける。」
「補題7.3：『strict linear derivation』が定義されていない。」

\Reply
投稿時第7節に相当する部分を改訂版第8節に短縮した。本文には定理8.3で使う単一spine正規形（命題8.1）と短い生成語の上界（補題8.2）を残し，補題8.2には直後に短い証明を置いた。付録Bでは統一的なspine鎖を使って正規化を証明した。未定義だった \emph{strict linear derivation} は削除した。
\Location{命題8.1と周辺の議論、補題8.2とその証明、付録B。}

\Comment{R1.32}
「命題7.2の証明に2ページも使う必要はない。全く同じ制限が定義されているかは知らないが、これは本質的でないことが知られているいくつかの制限の組合せであり、多くの読者は自明か民間伝承的事実として受け入れるだろう。」
「補題7.5と7.6はあまりにも自明である。」

\Reply
同意する。旧2ページの正規化証明はR1.31で説明した命題8.1（証明は付録B）と補題8.2（短い証明は本文）に整理した。旧補題7.5と7.6は削除するか各証明に吸収した。付録Bでは古典的正規化操作について Autebert--Berstel--Boasson (1997) と Hopcroft--Motwani--Ullman (2007) を引用し，二項規則のラッパー子がちょうど一つというspine構造に加え，文法サイズと構成時間の多項式上界を確認した。
\Location{命題8.1と周辺の議論、補題8.2とその証明、定理8.3とその後の補足、付録B。}

\subsection*{投稿時第8--10節}
\Comment{R1.33}
「アルファベット $\{d,u,;\}$ は適切でない。とくにセミコロンは本文中では読みにくく、$u$ は他の箇所で文字列を表すことがほとんどなので混乱の原因になる。……」
「定理8.6の証明は簡単にすべきだ。証拠語を短くするために複雑な式やギリシア文字を多用しても本質ではなく、読者の負担を増す。……」

\Reply
同意する。$d\mapsto\uparrow$（増加），$u\mapsto\downarrow$（減少），$;\mapsto\#$（リセット），$p\mapsto\rho$ とし，査読者が提案した証拠語の本質的な形を採用した。定理9.2の二語 $(\uparrow\downarrow)^k\uparrow^{\rho-1}\#^\ell$ と $(\uparrow\downarrow)^k\uparrow^\rho\#^\ell$ は空文脈と固定窓の端を共有するが，先行する $\uparrow$ で分布を区別できる。命題3.3直後の判定条件により長い議論を避けた。鋭い境界 $\rho\ge2$ は $\mathrm{CTR}_1$ が $(1,1)$-代入可能であることを脚注で説明した。
\Location{第9.2節。}

\Comment{R1.34}
「この話題のために小節を一つ丸ごと使う必要はない。系8.7が何を確立するのか不明である。……命題3.5と定理8.6には意味があるが、系8.7にまとめることで重要な貢献をかえって見えにくくしている。」
「『CCL の語のブロックとは、区切り記号 $;$ で区切られた部分語のそれぞれを指すことに注意する』というのは注意書きではなく新しい定義である。しかも定理8.6の証明はブロックという語を全く別の意味で使っている。」
「第8--10節はいずれも学習対象の言語クラスの豊かさに関するものであり、一つの節に統合し、説明と証明を簡潔にすれば論文の全体像が捉えやすくなる。」

\Reply
個別に学習可能な言語族の合併が学習可能とは限らないことに同意し，旧系8.7を削除した。$\RS$ は表現力の比較にだけ使い，学習可能性は固定 $h$ の文脈自由部分クラスにのみ主張する。旧命題3.5と定理8.6に対応する命題3.3と定理9.2は別の主張として残す。三つの境界関連の節は第9節に統合した。曖昧な \emph{block} の術語を因子と高さの記述に置き換え，$\mathrm{CTR}_\rho$ の型付け $h_\rho$ を第9.2節のDFAの遷移モノイド準同型として定めた。Yoshinaka (2008, Proposition 1) の正規言語の分離結果を明記した。投稿時第10節の $\Delta^*=(\{a^nb^n:n\ge0\})^*$ は命題9.1に短縮し（R1.35），不要なKleene星の例と $\mathrm{CCL}_1$ の付録は削除した。
\Location{第9節。}

\Comment{R1.35}
「定義9.1：記号 $\#$ を定義し、第9節を通じて一貫して使用すること。」
「命題9.2は自明で証明不要である。命題9.6も Ginsburg and Greibach, \emph{Deterministic context-free languages}, Information and Control 9 (1966), 620--648 に既に記載されている周知の事実である。」
「補題9.8：Łukasiewicz 言語と Dyck 言語の関係も周知である。」
「命題10.3の証明も省くか大幅に短縮できる。命題10.5の証明では定義3.3との記号の整合を検討されたい。」

\Reply
リセット記号 $\#$ は一度だけ定義して一貫して使う。$D_1$ と Łukasiewicz 言語の関係は Ginsburg--Greibach (1966) および Autebert--Berstel--Boasson (1997) を引用した。投稿時命題10.3の証明は命題9.1に短縮した。決定性文脈自由性はプッシュダウンによる簡潔な議論とし，代入可能性は入口高さ集合 $\mathcal Q(x)$ と脚注の連結恒等式により，$ba$ を含む因子と $a^*b^*$ の因子に分けて直接証明した。他の長い反例証明は命題3.3直後の判定条件または初等的なDFAの記述に縮約した。
\Location{第9節。}

\Comment{R1.36}
「補題9.11は $\sim_h$-代入可能性の基本的性質である。基本的な数学的性質を付録で集めるなら、そこに置いてもよいだろう。」

\Reply
投稿時の基本性質の付録は残さず，正規言語との共通部分・逆準同型像の結果を，消去因子を区別する非空因子の規約のもとで命題3.2にまとめた。残る右商の性質は直後で使う二行の補題9.6で述べ，一事実のための付録は設けない。
\Location{補題9.6。}

\subsection*{投稿時第11節}
\Comment{R1.37}
「$\RS\cap\mathrm{CFL}$ が Clark の合同文脈自由言語（congruential CFL）のクラスに含まれることにも触れる価値がある。」

\Reply
命題9.7として追加した。Clark (2010) は代入可能言語および $(k,\ell)$-代入可能言語を合同文脈自由言語族に含め、Dyck 言語の例も挙げている。改訂稿の命題は、固定 $h$ の各スライスについて包含関係を示し、したがって $\RS(\Gamma)\cap\mathrm{CFL}$ に拡張する。$\{a,b\}$ 上の真の包含は、合同的でありながら $\RS$ の外にある $D_1$ が示す。
\Location{命題9.7。}
